\documentclass[12pt,a4paper]{report}

\usepackage{amsmath, amssymb}
\usepackage{geometry}
\usepackage{setspace}
\usepackage{natbib}
\usepackage{hyperref}

\geometry{margin=1in}
\onehalfspacing

\title{Literature Review: Deep Learning for Ventilator Pressure Time-Series Modeling}
\author{}
\date{}

\begin{document}
	
	\maketitle
	
	\chapter*{Literature Review}
	
	\section*{Long Short-Term Memory Networks}
	
	Recurrent Neural Networks (RNNs) suffer from vanishing and exploding gradients when modeling long temporal dependencies. To address this limitation, Hochreiter and Schmidhuber (1997) introduced the Long Short-Term Memory (LSTM) architecture.
	
	The LSTM cell is defined as:
	
	\begin{align}
		f_t &= \sigma(W_f x_t + U_f h_{t-1} + b_f) \\
		i_t &= \sigma(W_i x_t + U_i h_{t-1} + b_i) \\
		\tilde{c}_t &= \tanh(W_c x_t + U_c h_{t-1} + b_c) \\
		c_t &= f_t \odot c_{t-1} + i_t \odot \tilde{c}_t \\
		o_t &= \sigma(W_o x_t + U_o h_{t-1} + b_o) \\
		h_t &= o_t \odot \tanh(c_t)
	\end{align}
	
	The additive update of the memory cell enables stable gradient flow across long sequences.
	
	In ventilator pressure modeling, each breath consists of 80 timesteps, where pressure depends on airflow, resistance, and compliance. LSTMs are therefore well-suited for capturing nonlinear temporal dependencies in respiratory dynamics.
	
	\section*{Sequence-to-Sequence Learning}
	
	Sutskever et al. (2014) introduced the sequence-to-sequence (Seq2Seq) framework using stacked LSTMs.
	
	The conditional probability is factorized as:
	
	\[
	P(Y|X) = \prod_{t=1}^{T'} P(y_t \mid y_{<t}, c)
	\]
	
	This formulation supports joint prediction of an entire sequence.
	
	In the ventilator setting:
	
	\[
	\{x_t\}_{t=1}^{80} \rightarrow \{p_t\}_{t=1}^{80}
	\]
	
	This motivates training the model to predict full pressure trajectories rather than independent timesteps.
	
	\section*{Deep Learning for Time Series}
	
	Fawaz et al. (2019) provide a comprehensive review of deep learning methods for time series. They demonstrate the effectiveness of deep architectures in modeling nonlinear multivariate signals.
	
	Key advantages include:
	
	\begin{itemize}
		\item End-to-end representation learning
		\item Nonlinear temporal modeling
		\item Scalability to high-dimensional inputs
	\end{itemize}
	
	Ventilator signals exhibit nonlinear airflow-pressure coupling, making LSTM-based approaches appropriate.
	
	\section*{Model Interpretability with SHAP}
	
	Lundberg and Lee (2017) introduced SHAP, a game-theoretic method for feature attribution.
	
	The Shapley value for feature $i$ is:
	
	\[
	\phi_i =
	\sum_{S \subseteq F \setminus \{i\}}
	\frac{|S|!(|F|-|S|-1)!}{|F|!}
	\left[
	f(S \cup \{i\}) - f(S)
	\right]
	\]
	
	SHAP guarantees local accuracy and consistency.
	
	In ventilator modeling, SHAP enables timestep-level analysis and identification of instability drivers in high-error sequences.
	
	\chapter*{Bibliography}
	
	\begin{thebibliography}{9}
		
		\bibitem{hochreiter1997lstm}
		Hochreiter, S., \& Schmidhuber, J. (1997).
		Long Short-Term Memory.
		\textit{Neural Computation}, 9(8), 1735–1780.
		
		\bibitem{sutskever2014seq2seq}
		Sutskever, I., Vinyals, O., \& Le, Q. V. (2014).
		Sequence to Sequence Learning with Neural Networks.
		\textit{Advances in Neural Information Processing Systems}.
		
		\bibitem{fawaz2019dlts}
		Fawaz, H. I., Forestier, G., Weber, J., Idoumghar, L., \& Muller, P.-A. (2019).
		Deep Learning for Time Series Classification: A Review.
		\textit{Data Mining and Knowledge Discovery}.
		
		\bibitem{lundberg2017shap}
		Lundberg, S. M., \& Lee, S.-I. (2017).
		A Unified Approach to Interpreting Model Predictions.
		\textit{Advances in Neural Information Processing Systems}.
		
	\end{thebibliography}
	
\end{document}